% AI4MATH 2026 workshop submission.
% Title: "pythia: a Lean 4 tactic library for statistical proof automation"
%
% NOTE on template: AI4MATH 2026 (https://ai4math2026.github.io/#cfp) does
% not publish a venue-specific style file. AI4MATH workshops in prior
% cycles inherited the host conference's template (NeurIPS 2024, ICML
% 2025). Aidan's directive (2026-04-26) selects ICML formatting since
% AI4MATH 2026 sits in the ML-conference orbit. We vendor the official
% ICML 2026 styles (icml2026.sty / icml2026.bst, dated 2025-10-29).
%
% Compile: bash build.sh  (runs pdflatex x3 + bibtex).

\documentclass{article}

% Style files vendored under styles/. TEXINPUTS in build.sh adds
% styles/ to the search path so \usepackage{icml2026} resolves.
\usepackage{icml2026}

\usepackage{microtype}
\usepackage{graphicx}
\usepackage{booktabs}
\usepackage{amsmath,amssymb,amsthm,mathtools}
\usepackage{newunicodechar}
% Map math-Unicode chars that appear in verbatim Lean blocks to LaTeX
% equivalents so pdflatex doesn't choke on them.
\newunicodechar{ℕ}{\ensuremath{\mathbb{N}}}
\newunicodechar{ℝ}{\ensuremath{\mathbb{R}}}
\newunicodechar{ℤ}{\ensuremath{\mathbb{Z}}}
\newunicodechar{∀}{\ensuremath{\forall}}
\newunicodechar{∃}{\ensuremath{\exists}}
\newunicodechar{≤}{\ensuremath{\leq}}
\newunicodechar{≥}{\ensuremath{\geq}}
\newunicodechar{≠}{\ensuremath{\neq}}
\newunicodechar{∑}{\ensuremath{\sum}}
\newunicodechar{∞}{\ensuremath{\infty}}
\newunicodechar{∈}{\ensuremath{\in}}
\newunicodechar{↦}{\ensuremath{\mapsto}}
% newunicodechar's built-in defaults already map →, ←, ·, ⟨, ⟩ to the
% LaTeX equivalents, so we don't redefine them — that would produce
% "Redefining Unicode character" warnings on every build pass.
\newunicodechar{α}{\ensuremath{\alpha}}
\newunicodechar{β}{\ensuremath{\beta}}
\newunicodechar{γ}{\ensuremath{\gamma}}
\newunicodechar{δ}{\ensuremath{\delta}}
\newunicodechar{ε}{\ensuremath{\varepsilon}}
\newunicodechar{σ}{\ensuremath{\sigma}}
\newunicodechar{τ}{\ensuremath{\tau}}
\newunicodechar{λ}{\ensuremath{\lambda}}

\usepackage{hyperref}
% \theHalgorithm is auto-defined by hyperref+algorithm in some
% distributions and not others. \providecommand is the portable form.
\providecommand{\theHalgorithm}{\arabic{algorithm}}
\usepackage[capitalize,noabbrev]{cleveref}
\usepackage{listings}
\usepackage{xcolor}
\usepackage{tikz}
\usetikzlibrary{positioning,arrows.meta}
\usepackage{enumitem}
\usepackage{float}
\usepackage[section]{placeins}
\setlength{\emergencystretch}{6em}

\theoremstyle{plain}
\newtheorem{theorem}{Theorem}[section]
\newtheorem{proposition}[theorem]{Proposition}
\newtheorem{lemma}[theorem]{Lemma}
\newtheorem{corollary}[theorem]{Corollary}
\theoremstyle{definition}
\newtheorem{definition}[theorem]{Definition}
\newtheorem{assumption}[theorem]{Assumption}
\theoremstyle{remark}
\newtheorem{remark}[theorem]{Remark}

\newcommand{\R}{\mathbb{R}}
\newcommand{\N}{\mathbb{N}}
\newcommand{\E}{\mathbb{E}}
\newcommand{\Prob}{\mathbb{P}}

% Listings setup. We use plain verbatim for Lean fragments rather than
% a custom lstdefinelanguage; the paper is about the tactic surface,
% not Lean source archaeology, so syntax highlighting is not needed.
% (Avoiding lstdefinelanguage also keeps the preamble paper_lint clean
% on '--' Lean-comment lexer hints.)
\lstset{
  basicstyle=\ttfamily\small,
  breaklines=true,
  showstringspaces=false,
  columns=fullflexible,
}

\icmltitlerunning{pythia: a Lean 4 tactic library for statistical proof automation}

\begin{document}

\twocolumn[
\icmltitle{\texorpdfstring{pythia: a Lean 4 tactic library for \\
           statistical proof automation}{pythia: a Lean 4 tactic library for statistical proof automation}}

\icmlsetsymbol{equal}{*}

\begin{icmlauthorlist}
\icmlauthor{Anonymous Authors}{anon}
\end{icmlauthorlist}

\icmlaffiliation{anon}{Affiliation redacted for double-blind review.}

\icmlcorrespondingauthor{Anonymous Authors}{anon@example.org}

\icmlkeywords{Lean 4, proof automation, statistics, anytime-valid inference, tactic library, automated theorem proving}

\vskip 0.3in
]

\printAffiliationsAndNotice{}

% Abstract: AI4MATH 2026 workshop submission.
\begin{abstract}
Lean 4 with Mathlib ships strong general automation (\texttt{aesop},
\texttt{simp}, \texttt{linarith}, \texttt{polyrith}, \texttt{nlinarith},
\texttt{positivity}, \texttt{measurability}, \texttt{bound},
\texttt{gcongr}). It does not ship a hammer for applied
mathematics. The moment a goal mentions a martingale, a stopping
time, a production function, an ordinary differential equation
(ODE) Lyapunov candidate, a Wasserstein distance, or a quantum
density operator, the generic tactics have nothing to apply.
We present \texttt{pythia}, a Lean 4 tactic library for
applied-mathematics proof automation. The headline tactic
\texttt{pythia!} orchestrates a nine-rung deterministic ladder.
The rungs are \texttt{stat\_simp}, the
\texttt{linarith}/\texttt{nlinarith}/\texttt{polyrith} chain,
\texttt{positivity}, \texttt{aesop} on the registered
\texttt{Pythia} ruleset, the \texttt{@[stat\_lemma]} cascade,
\texttt{z3\_check}, \texttt{cvc5\_check}, \texttt{vampire\_check}
or \texttt{e\_check}, and \texttt{disprove}. Each rung tries
fail-fast with a bounded budget. The verbose variant
\texttt{pythia?} reports the closing rung and per-rung wall-clock
timing for diagnostic use. The library ships ten registered
tactics, attribute-driven registries (\texttt{@[stat\_lemma]},
\texttt{@[stats\_ineq]}, \texttt{@[prob\_simp]},
\texttt{@[stat\_simp]}, \texttt{@[anytime\_valid\_lemma]}, plus
the per-domain calculator typeclass), and a fifteen-domain proof
library with thirty-two cross-domain theorems carrying matched
Python sim runners (property-based testing, deterministic sweeps,
mutation testing). Domains span probability and statistics,
biology, chemistry, economics, control theory, optimal transport,
quantum information, stochastic calculus, game theory,
thermodynamics, numerical analysis, queueing, and operations
research. Every public theorem is axiom-clean against the standard
Lean kernel set ($\texttt{propext}$, $\texttt{Classical.choice}$,
$\texttt{Quot.sound}$). The cross-prover ladder is offline-first
and free of large language models (LLMs): \texttt{z3\_check},
\texttt{cvc5\_check}, \texttt{vampire\_check}, and \texttt{e\_check}
dispatch to local Satisfiability Modulo Theories (SMT) and
Automated Theorem Prover (ATP) binaries with kernel-checked
reconstruction in the CoqHammer style of
\citet{czajkakaliszyk2018coqhammer}. LLM-augmented closure lives
in a separate companion software development kit (SDK) that the
library does not depend on. Hard structural Mathlib gaps are
optionally routed to the Aristotle automated theorem prover
\citep{harmonic2025aristotle} for development-time restocking,
with audit gates against vacuous-truth closures. The library is
pinned to Lean 4.28 + Mathlib v4.28.0 and released open-source.
\end{abstract}

\section{Introduction}
\label{sec:intro}

Lean 4 with Mathlib ships strong general-purpose proof automation:
\texttt{aesop} for forward search, \texttt{simp} for rewriting,
\texttt{linarith} and \texttt{polyrith} for linear and polynomial
arithmetic, \texttt{nlinarith} for non-linear arithmetic,
\texttt{positivity} for sign goals, \texttt{measurability} for the
$\sigma$-algebra fragment, \texttt{bound} for monotone composition,
\texttt{gcongr} for generalised congruence
\citep{moura2021lean, mathlib2020community, mathlib2025}.
Applied-mathematics proofs route through these general tactics,
typically by hand decomposition into a chain of \texttt{simp} sets,
\texttt{linarith} hypotheses, and Mathlib lemma selection. None of
the general tactics know about each other. The user is the
orchestrator. There is no domain hammer for applied mathematics.

CoqHammer \citep{czajkakaliszyk2018coqhammer} dispatches Coq goals
to external automated theorem provers (ATPs) such as Vampire, E,
and CVC4, with kernel-checked reconstruction. Sledgehammer
\citep{paulsonblanchette2010threeyears, blanchettesledgehammer2011}
plays the same role in Isabelle/Higher Order Logic (HOL). Lean 4 has no comparable
hammer. This paper does not attempt the fully-general case across
all of Lean. We address the applied-mathematics surface: a
registered library of concentration, anytime-valid, optimal
transport, control-theoretic, and information-theoretic lemmas
plus a fan-out across local Satisfiability Modulo Theories (SMT)
and ATP solvers, packaged as a single deterministic tactic call.

\paragraph{Contribution.}
We present \texttt{pythia}, a Lean 4 tactic library for
applied-mathematics proof automation. The headline tactic
\texttt{pythia!} orchestrates a nine-rung deterministic ladder
under one call. The first five rungs are pure-Lean tactics. They are
\texttt{stat\_simp} (curated probability and ENNReal normal
forms), the \texttt{linarith}/\texttt{nlinarith}/\texttt{polyrith}
chain, \texttt{positivity}, \texttt{aesop} on the registered
\texttt{Pythia} ruleset, and the \texttt{@[stat\_lemma]}
shape-dispatch cascade. The next three rungs are external solvers
with kernel-checked reconstruction. They are \texttt{z3\_check}
(linear real arithmetic via Z3), \texttt{cvc5\_check} (bit-vector
and linear real arithmetic via CVC5), and \texttt{vampire\_check}
or \texttt{e\_check} (first-order logic, FOL, via Vampire or E).
The final rung is \texttt{disprove}, a counterexample finder via
Z3 model extraction. Each rung tries fail-fast with a per-rung
budget, and first-to-close wins. The
verbose variant \texttt{pythia?} reports the closing rung and
per-rung wall-clock timing for diagnostic use. \texttt{pythia}
ships ten registered tactics in total. Six attribute-driven
registries cover lemma routing: \texttt{@[stat\_lemma]},
\texttt{@[stats\_ineq]}, \texttt{@[prob\_simp]},
\texttt{@[stat\_simp]}, \texttt{@[anytime\_valid\_lemma]}, plus
the \texttt{DomainCalculator} typeclass for per-domain numerical
calculators. The proof library spans fifteen domains and carries
thirty-two cross-domain theorems with matched Python sim runners.
Sim runners exercise property-based testing, deterministic
sweeps, and mutation testing.
The shipped library is axiom-clean against
$\{\texttt{propext}, \texttt{Classical.choice}, \texttt{Quot.sound}\}$
and pinned to Lean 4.28 + Mathlib v4.28.0. The cross-prover ladder
is offline-first and free of large language models (LLMs). SMT and
ATP backends are deterministic decision procedures with verifiable
Lean reconstruction. LLM-augmented closure lives in a separate
companion software development kit (SDK) that the library does not
depend on. Hard structural Mathlib gaps are optionally routed to
the Aristotle automated theorem prover
\citep{harmonic2025aristotle} for development-time restocking only,
with kernel-checked reconstruction inside Lean and an audit gate
against vacuous closures.

\paragraph{Case study (vacuous-truth catch).}
On 2026-04-26 the Aristotle prover closed a candidate theorem in
our \texttt{wald\_wolfowitz\_optimal} development by routing
through a contradictory premise. A Bool partition combined with a
log-boundary constraint forced both $\alpha \geq 1$ and $\alpha <
1$, closing the goal vacuously. The closure was kernel-checked.
The theorem statement was the bug. The library audit caught the
issue and forced a restatement to the two-measure form. We treat
this kind of catch as a first-class library responsibility, not as
a separate review pass.

\paragraph{Paper organisation.}
Section~\ref{sec:tactic-surface} describes the ten-tactic surface
and the registry attributes.
Section~\ref{sec:library} groups the shipped library content by
domain. Section~\ref{sec:aristotle} describes the Aristotle
integration as a development-time oracle and the vacuous-truth
refusal policy.
Section~\ref{sec:discussion} compares against aesop, Sledgehammer,
CoqHammer, and Mathlib's existing tactic suite, states limitations,
and reports a small head-to-head between Aristotle, the DeepSeek
Prover-V2 (DSPv2) model, and Claude Opus 4.6 on a subset of
FormalAVS \citep{anonymous2026formalavs}. Full numbers are in the
dataset paper.

\paragraph{Reproduction and code.}
The \texttt{pythia} library is open source. The URL is redacted
for double-blind review. Axiom audit transcripts and reproduction
recipes are listed in Appendix~\ref{sec:appendix-repro}. The
dataset of theorem-proving benchmarks (FormalAVS,
\citet{anonymous2026formalavs}) is reported in a separate paper.

\section{The pythia tactic surface}
\label{sec:tactic-surface}

\texttt{pythia} ships five user-facing tactics, one named aesop
ruleset, and three attribute-driven registries. The design follows
the CoqHammer pattern \citep{czajkakaliszyk2018coqhammer}: a
domain-specific dispatcher with extensible attribute-driven lemma
selection, a fall-through to general-purpose automation, and
kernel-checked reconstruction for any external-oracle steps.

\subsection{Five tactics}
\label{sec:tactic-surface-five}

\begin{table}[h]
\centering
\small
\begin{tabular}{lp{4.5cm}}
\toprule
\textbf{Tactic} & \textbf{Role} \\
\midrule
\texttt{pythia} & Top-level dispatcher. Tries registered \texttt{stat\_lemma}s, then aesop ruleset \texttt{Pythia}, then \texttt{linarith}. \\
\texttt{stats\_ineq} & Discharges concentration and tail-bound inequalities (Hoeffding, Bernstein, sub-Gaussian, sub-gamma). \\
\texttt{prob\_simp} & Probability normalization rewrites. Pushes integrals, splits indicator products, reduces to MGF form. \\
\texttt{anytime\_valid} & Closes Ville-bound goals for anytime-valid confidence sequences. Specializes to Howard-Ramdas, betting, vector. \\
\texttt{z3\_check} & Linear-real-arithmetic Z3 oracle with linarith reconstruction. Kernel-checked. \\
\bottomrule
\end{tabular}
\caption{The five \texttt{pythia} tactics. \texttt{pythia} is the
default entry point. The four specialized tactics are exposed for
direct invocation when the user knows the goal shape.}
\label{tab:tactic-surface}
\end{table}

% TODO: each tactic deserves a one-paragraph design note. asabi to draft.

\paragraph{\texttt{pythia} (top-level dispatcher).}
The default entry point. Tries the registered statistical lemma
database, then the named aesop ruleset \texttt{Pythia}, then
\texttt{linarith}. The fall-through ordering is fixed. The lemma
database and aesop ruleset are user-extensible.

\paragraph{\texttt{stats\_ineq}.}
Discharges sub-Gaussian, sub-gamma, Bernstein-Bennett, and
Hoeffding-shape inequalities. The lemma set is registered via
\texttt{@[stats\_ineq]}. Reconstruction is internal Lean and
axiom-clean.

\paragraph{\texttt{prob\_simp}.}
Probability-domain rewrite set. Pushes integrals through linear
operators, splits indicator products, and reduces martingale
expressions to MGF form for downstream use by
\texttt{stats\_ineq}. The rewrite set is registered via
\texttt{@[prob\_simp]}.

\paragraph{\texttt{anytime\_valid}.}
Closes Ville-type maximal-inequality goals. Specializes to the
Howard-Ramdas \citep{howard2021time}, betting
\citep{waudbysmith2024estimating}, and vector
\citep{whitehouse2025timeuniform} families, dispatching by goal
shape. The case-split table is small (currently four families) and
hard-coded. Future families register via \texttt{@[anytime\_valid]}.

\paragraph{\texttt{z3\_check}.}
External Z3 oracle for linear real arithmetic, with
\texttt{linarith} reconstruction inside Lean. Z3 supplies the proof
witness. \texttt{linarith} reconstructs the witness step inside
the kernel. The kernel is the final word. Z3 cannot smuggle axioms.
This follows the CoqHammer reconstruction pattern
\citep{czajkakaliszyk2018coqhammer}.

\subsection{Registry attributes}
\label{sec:tactic-surface-registries}

The library uses three attribute-driven registries:

\begin{itemize}[leftmargin=1.3em,itemsep=0.2em]
\item \texttt{@[stat\_lemma]}: lemma database for the top-level
  \texttt{pythia} dispatcher.
\item \texttt{@[stats\_ineq]}: tail-bound and concentration
  lemmas dispatched by \texttt{stats\_ineq}.
\item \texttt{@[prob\_simp]}: probability rewrite rules dispatched
  by \texttt{prob\_simp}.
\end{itemize}

User extensions consist of tagging library lemmas with the
appropriate attribute. No fork is required.

\subsection{The \texttt{Pythia} aesop ruleset}
\label{sec:tactic-surface-aesop}

The library declares a named aesop ruleset \texttt{Pythia} containing
the registered \texttt{stat\_lemma} entries. Direct invocation is
supported via \texttt{aesop (rule\_sets := [Pythia])}, equivalent to
the \texttt{pythia} dispatcher's middle stage. This gives users a
finer-grained interface when the dispatcher's outer wrapping is not
desired.

\subsection{Architecture}
\label{sec:tactic-surface-arch}

% TODO: insert TikZ figure. Sketch:
%   goal -> [registry lookup, three attribute-driven sets]
%        -> [aesop (rule_sets := [Pythia])]
%        -> [linarith fall-through]
%   side-channel: z3_check (separate user-invoked tactic, not
%   part of the pythia dispatcher chain).
\begin{figure}[h]
\centering
\framebox[0.95\columnwidth]{\parbox{0.85\columnwidth}{
\textbf{Architecture (TODO: replace with TikZ).}
The \texttt{pythia} dispatcher tries the
\texttt{@[stat\_lemma]} database first, then the
\texttt{Pythia} aesop ruleset, then \texttt{linarith}. Each stage
is independently Lean-kernel-checked. \texttt{z3\_check} is a
separate user-invoked tactic with its own
linarith-based reconstruction step, not part of the dispatcher
chain.
}}
\caption{The \texttt{pythia} dispatcher chain. Three stages, each
kernel-checked independently. \texttt{z3\_check} sits outside the
chain.}
\label{fig:arch}
\end{figure}

\paragraph{Why this shape.}
The three-stage chain mirrors the empirical structure of statistical
proofs in our library: many goals close on a single lemma application
(stage 1), some require forward search through a lemma set (stage 2),
some reduce to a linear arithmetic obligation after the rewrite
(stage 3). Routing by goal shape is not necessary at the top level.
The cost of trying stages in order is small relative to the cost of
incorrect early termination.

% TODO: add brief usage example. asabi to draft 5-10 line code block
% showing pythia closing a Hoeffding instance.

\section{Library content}
\label{sec:library}

The shipped \texttt{pythia} library spans a tactic layer plus a
proof-content layer of fifteen applied-math domains. The full
manifest at the time of writing carries thirty-two cross-domain
theorems with matched Python sim runners (property-based testing,
parameter sweeps, mutation testing) plus four hundred and sixty
total theorem and lemma declarations across the
\texttt{Pythia/} source tree, fifty-five of which are tagged
\texttt{@[stat\_lemma]} and registered into the headline cascade.
Every theorem listed here is axiom-clean against
$\{\texttt{propext}, \texttt{Classical.choice}, \texttt{Quot.sound}\}$
under the Lean 4.28 + Mathlib v4.28.0 toolchain pin. We summarise
by group. The appendix lists modules in detail.

\subsection{Anytime-valid inference}
\label{sec:library-av}

Howard-Ramdas confidence sequences \citep{howard2021time}, betting-CS
\citep{waudbysmith2024estimating}, and the vector-valued extension
\citep{whitehouse2025timeuniform} are each backed by a Ville-bound
specialization in the \texttt{anytime\_valid} tactic. The
\texttt{Pythia.HowardRamdasCS} and \texttt{Pythia.BettingCS} modules
expose the admissibility certificates of the e-process formalism of
\citet{ramdasgrunwaldvovkshafer2023}. The supporting Ville machinery
lives in \texttt{Pythia.VilleSupermartingale} (countable-time and
finite-horizon variants).

\subsection{Concentration and tail bounds}
\label{sec:library-conc}

\texttt{stats\_ineq} dispatches sub-Gaussian and sub-gamma tail
bounds, the Bernstein-Bennett MGF bridge, and Hoeffding-style
inequalities for bounded random variables. Modules:
\texttt{Pythia.SubGaussianMG} (sub-Gaussian martingale concentration),
\texttt{Pythia.SubGamma} (sub-gamma tail bound),
\texttt{Pythia.MGFBoundedSubGamma} (the bounded-RV Bernstein-Bennett
MGF bound, $\mathbb{E}[\exp(\lambda X)] \leq \exp(\sigma^2 \lambda^2 /
(2(1 - b\lambda/3)))$ for centred $|X| \leq b$), and
\texttt{Pythia.Bernstein} (the iid and martingale Bernstein
applications). The MGF bridge composes with \texttt{prob\_simp} for
the integral side and with general Mathlib infrastructure
(\texttt{measurability}, \texttt{integrability}) for side conditions.

\subsection{E-detectors and merging}
\label{sec:library-edet}

The e-process formalism of
\citet{ramdasgrunwaldvovkshafer2023, shinramdasrinaldo2024}
yields uniform-time validity by combining e-processes via averaging.
\texttt{Pythia.EDetector} declares \texttt{EProcess} as a non-negative
process bounded in expectation under the null and \texttt{EDetector}
as the first-passage stopping rule with anytime-valid Type-I error
$\leq \alpha$. The \texttt{combine\_eprocesses\_avg} lemma (average
merge) is the workhorse. Product-merge variants register alongside.
The average-merge case is axiom-clean.

\subsection{Matrix concentration}
\label{sec:library-matrix}

Matrix Bernstein \citep{tropp2012usermatrix} and the Lieb concavity
inequality \citep{lieb1973} are the two matrix-side modules.
\texttt{Pythia.MatrixLieb} states the Lieb concavity step and the
operator-monotone log composition.
\texttt{Pythia.MatrixBernstein} carries the iid matrix-Bernstein head
theorem. \texttt{Pythia.MatrixBernsteinFull} carries the
matrix-martingale extension. The dependent operator-trace inequalities
route through Aristotle as oracle (Section~\ref{sec:aristotle}).

\subsection{Cross-domain coverage (15 fields)}
\label{sec:library-other}

The cross-domain manifest covers fifteen fields with Lean proofs
and matched Python sim runners. Each entry is a closed-form
theorem with property-based-test mutation coverage in the
companion runner.

\begin{itemize}[leftmargin=1.3em,itemsep=0.2em]
\item \textbf{Probability and statistics:} the anytime-valid /
  concentration / matrix sections above plus \texttt{Pythia.Risk.}
  \texttt{CVaR}~\citep{rockafellaruryasev2000} and
  \texttt{Pythia.Asymptotics.} \texttt{DeltaMethod} (scalar) /
  \texttt{DeltaMethodMulti} (multivariate).
\item \textbf{Information theory:} \texttt{Pythia.InfoTheory.}
  \texttt{BretagnolleHuberBinary} (binary-alphabet bound),
  \texttt{Pythia.InfoTheory.DataProcessing} (data-processing
  inequality), \texttt{Pythia.InfoTheory.BinaryEntropy} (Shannon
  binary entropy).
\item \textbf{Stochastic calculus:}
  \texttt{Pythia.Stochastic.ItoIsometryFiniteDim} (finite-dim Itô
  isometry), \texttt{Pythia.BDG} (Burkholder-Davis-Gundy at
  $p = 2$), \texttt{Pythia.MeasureTheory.}
  \texttt{ConditionalJensen}, \texttt{Pythia.StochasticApproximation.}
  \texttt{RobbinsSiegmund} / \texttt{RobbinsMonro} /
  \texttt{Dvoretzky}, \texttt{Pythia.TimeSeries.WoldDecomposition}
  \citep{wold1938}, \texttt{Pythia.TimeSeries.NeweyWest}.
\item \textbf{Quantum information:}
  \texttt{Pythia.Quantum.VonNeumannEntropyNonnegTwoState}
  (qubit von Neumann entropy non-negativity).
\item \textbf{Optimal transport:}
  \texttt{Pythia.OptimalTransport.WassersteinDistanceNonneg}
  (discrete L1 Wasserstein non-negativity).
\item \textbf{Game theory:}
  \texttt{Pythia.GameTheory.MinimaxTwoStrategyBound}
  (von Neumann two-strategy minimax inequality).
\item \textbf{Biology:} \texttt{Pythia.Bio.Population}
  (Hardy-Weinberg, Lotka-Volterra equilibria, SIR conservation),
  \texttt{Pythia.Bio.MichaelisMentenSaturation} (enzyme
  saturation), \texttt{Pythia.Bio.MassAction} (CRN ODE
  conservation), \texttt{Pythia.Bio.Phylogenetics}.
\item \textbf{Chemistry:} \texttt{Pythia.Chemistry.Arrhenius}
  (rate-law positivity), \texttt{Pythia.Chemistry.}
  \texttt{HendersonHasselbalch} (pH monotonicity),
  \texttt{Pythia.Chemistry.MassActionConservation}.
\item \textbf{Economics:} \texttt{Pythia.Economics.CobbDouglas}
  (constant returns to scale), \texttt{Pythia.Economics.CRRA}
  (utility concavity), \texttt{Pythia.Economics.CAPM}
  (linear pricing), \texttt{Pythia.Economics.RiskNeutralCall}
  (non-negativity), \texttt{Pythia.Economics.Walras}
  (excess-demand identity).
\item \textbf{Engineering:}
  \texttt{Pythia.Engineering.RCTimeConstant},
  \texttt{Pythia.Engineering.SignalEnergy},
  \texttt{Pythia.Engineering.PowerDissipation}.
\item \textbf{Mechanical:}
  \texttt{Pythia.Mechanical.HookeSpring}
  (spring-potential non-negativity).
\item \textbf{Control theory:}
  \texttt{Pythia.Control.Lyapunov} (scalar),
  \texttt{Pythia.Control.LyapunovDiscrete} (discrete-time),
  \texttt{Pythia.Control.LyapunovODE}.
\item \textbf{Operations research / queueing:}
  \texttt{Pythia.Queueing.ErlangB},
  \texttt{Pythia.Queueing.LittlesLaw},
  \texttt{Pythia.OR.LittlesLaw}.
\item \textbf{Numerical analysis:}
  \texttt{Pythia.Numerical.NewtonQuadraticIterPos},
  \texttt{Pythia.Numerical.PicardLindelof},
  \texttt{Pythia.Numerical.Lyapunov},
  \texttt{Pythia.Numerical.KKT},
  \texttt{Pythia.Numerical.Kahan}.
\item \textbf{Thermodynamics:}
  \texttt{Pythia.Thermodynamics.CarnotEfficiencyUpperBound}.
\item \textbf{Mathlib retags + actuarial:}
  \texttt{Pythia.MathlibTags} (AM-GM, Markov, Cauchy-Schwarz with
  empirical companions), \texttt{Pythia.Actuarial.\{Pareto,
  Weibull, LogNormal\}}.
\end{itemize}

\subsection{Five highlights}
\label{sec:library-highlights}

We name five representative results that exercise the tactic surface.
Module-level statements appear in
Appendix~\ref{sec:appendix-inventory}.

\begin{enumerate}[leftmargin=1.3em,itemsep=0.2em]
\item \textbf{Howard-Ramdas Ville bound.} The maximal inequality for
  the self-normalized confidence sequence
  \citep{howard2021time}, in
  \texttt{Pythia.HowardRamdasCS}. Closes via \texttt{anytime\_valid}.
\item \textbf{Betting-CS coverage.} The supermartingale coverage
  certificate of \citet{waudbysmith2024estimating}, in
  \texttt{Pythia.BettingCS}. Axiom-clean closure.
\item \textbf{Bernstein-Bennett MGF bridge.} The bounded-RV MGF bound
  in \texttt{Pythia.MGFBoundedSubGamma}, routing between the
  variance-aware sub-gamma tail and the Bennett MGF form.
\item \textbf{combine\_eprocesses\_avg.} Average-merge of e-processes
  \citep{ramdasgrunwaldvovkshafer2023, shinramdasrinaldo2024} in
  \texttt{Pythia.EDetector}. Axiom-clean closure with full
  reconstruction.
\item \textbf{Bretagnolle-Huber binary.} The binary-case
  total-variation lower bound in
  \texttt{Pythia.InfoTheory.} \texttt{BretagnolleHuberBinary}. Standard in
  information-theoretic lower-bound arguments.
\end{enumerate}

\paragraph{Pinning.}
The library is pinned to Lean 4.28 + Mathlib v4.28.0. The pin is
necessary: the modules consume Mathlib's measure theory and
probability subdirectories, which evolve quickly. We track Mathlib
release tags rather than \texttt{master} for reproducibility.

\section{Aristotle integration as oracle}
\label{sec:aristotle}

A subset of statistical proof goals exceeds what aesop, simp,
linarith, and the registered \texttt{pythia} lemma sets close in
practice. We route these to the Aristotle automated theorem prover
\citep{harmonic2025aristotle} as an oracle. Aristotle returns a
Lean~4 proof script, the script is checked by the Lean kernel inside
the calling project, and the result is audited before integration.
Aristotle is one of two oracles in the library. The other is Z3
(Section~\ref{sec:tactic-surface-ten}).

\subsection{Submission pattern}
\label{sec:aristotle-pattern}

Each Aristotle call is a project tarball submission. The flow is:
state the theorem in Lean with a placeholder, package the project
without the Lean build cache (\texttt{.lake} excluded), submit via
the Aristotle endpoint, monitor for completion, integrate the
returned proof, and run the axiom audit. The audit is the
verification step that distinguishes a kernel-checked Aristotle
proof from a vacuous closure.

\subsection{Vacuous-truth refusal}
\label{sec:aristotle-vacuous}

ATPs can close a theorem by deriving a contradiction from its
hypotheses, then concluding anything. The kernel accepts such
proofs. The proof is sound. The theorem statement is the bug.

\paragraph{Case (2026-04-26).}
We submitted a candidate
\texttt{wald\_wolfowitz\_optimal} theorem. Aristotle returned a
closure that combined a Bool partition with a log-boundary constraint
to force both $\alpha \geq 1$ and $\alpha < 1$. The kernel accepted
the proof. The audit flagged the hypothesis chain as
self-contradictory. We rejected the closure, restated the theorem in
two-measure form (separating the null and alternative measures
explicitly), and resubmitted. The two-measure form was the intended
target. The original Bool-partition form was a statement bug.

\paragraph{Audit policy.}
The library treats vacuous closures as theorem-statement bugs, not as
prover successes. Concretely: every Aristotle-returned proof is
audited for contradictory premises before integration. Closures
flagged as vacuous block the corresponding theorem from the
axiom-clean library until the statement is reformed. This policy is
not unique to Aristotle. The same audit applies to any external
prover, including human proofs that close via a hypothesis
contradiction the author did not notice.

\subsection{Z3 as a second oracle}
\label{sec:aristotle-z3}

The \texttt{z3\_check} tactic dispatches linear-real-arithmetic goals
to Z3 \citep{demoura2008z3}. Z3 returns a verdict (unsat of the
negated goal). The verdict is then reconstructed inside Lean by
\texttt{linarith}. The Lean kernel is the final word. The pattern is
the CoqHammer \citep{czajkakaliszyk2018coqhammer} reconstruction
strategy adapted to Lean 4. Z3 cannot smuggle axioms.
\texttt{linarith} is axiom-clean and the reconstruction step
certifies the witness inside the kernel. If \texttt{linarith} cannot
reconstruct, the \texttt{z3\_check} call fails loud rather than
silently trusting Z3.

\subsection{When to call Aristotle}
\label{sec:aristotle-when}

The library convention is to attempt local closure first. If the
goal closes by \texttt{rfl}, \texttt{simp}, \texttt{positivity},
\texttt{linarith}, \texttt{nlinarith}, \texttt{polyrith}, or the
registered \texttt{stat\_lemma} database, no Aristotle call is made.
Aristotle is reserved for hard structural Mathlib gaps:
measure-theoretic $\sigma$-additivity steps, dominated-convergence
side conditions, matrix-concentration intermediate lemmas, and
similar obligations for which a hand-Mathlib closure either does not
exist or would require an order of magnitude more dispatch. The
vacuous-closure audit is the gate on integration. The dispatch
threshold is only the call decision.

\section{Comparison and discussion}
\label{sec:discussion}

\subsection{Comparison to existing automation}
\label{sec:discussion-comparison}

\paragraph{aesop.}
Aesop \citep{aesop2023} is the workhorse forward-search tactic in
Lean 4. \texttt{pythia} declares a named aesop ruleset and uses
aesop in stage 2 of the dispatcher chain. \texttt{pythia} is not a
replacement for aesop. It is a statistics-domain registry plus a
Z3 oracle plus a vacuous-closure audit, layered on top of aesop.

\paragraph{Sledgehammer.}
Sledgehammer \citep{paulsonblanchette2010threeyears,
blanchettesledgehammer2011} dispatches Isabelle/HOL goals to external
ATPs (Vampire, E, Z3, CVC4) and reconstructs the proof. There is no
direct Lean 4 equivalent. \texttt{pythia} addresses the
statistics-domain case via \texttt{z3\_check} for linear-real-arithmetic and
Aristotle as a structural oracle. A general-purpose Lean
Sledgehammer is out of scope for this paper.

\paragraph{CoqHammer.}
CoqHammer \citep{czajkakaliszyk2018coqhammer} dispatches Coq goals
to external ATPs with kernel-checked reconstruction. \texttt{z3\_check}
in \texttt{pythia} adopts the same reconstruction template
(witness from external prover, in-kernel checking via
\texttt{linarith}). CoqHammer is the design template. We do not
claim a Lean-side reimplementation of CoqHammer's full premise
selection.

\paragraph{Mathlib's existing tactic suite.}
Mathlib supplies \texttt{linarith}, \texttt{polyrith},
\texttt{nlinarith}, \texttt{positivity}, \texttt{measurability},
\texttt{bound}, \texttt{gcongr}. None of these is a statistics-domain
hammer. \texttt{pythia} composes with all of them. The fall-through
to \texttt{linarith} in stage 3 of the dispatcher chain is the
explicit composition.

\subsection{Aristotle vs DSPv2 vs Opus 4.6 (illustrative)}
\label{sec:discussion-comparison-empirical}

% TODO: pull 3-4 problems from FormalAVS where Aristotle closes
% and hand-Mathlib alone does not. asabi to populate from
% neurips-2026-dataset-paper tables.

Table~\ref{tab:illustrative-comparison} reports a small subset of
FormalAVS problems \citep{anonymous2026formalavs} on which
Aristotle, DSPv2 (a fine-tuned local-GPU drafter), and Opus 4.6
(an API model) were each given a single attempt. The full benchmark
numbers are in the dataset paper. The point of this subset is to
illustrate the regime in which Aristotle as oracle closes problems
that hand-Mathlib alone does not. Per-problem solve outcomes (close
or open) are kernel-checked.

\begin{table}[h]
\centering
\small
\begin{tabular}{lccc}
\toprule
\textbf{Problem} & \textbf{Aristotle} & \textbf{DSPv2} & \textbf{Opus 4.6} \\
\midrule
% TODO: replace placeholders with real solve/no-solve outcomes
% from the FormalAVS subset chosen for this illustrative comparison.
combine\_eprocesses\_avg & TODO & TODO & TODO \\
ville\_bound\_betting & TODO & TODO & TODO \\
matrix\_bernstein\_step & TODO & TODO & TODO \\
bretagnolle\_huber\_binary & TODO & TODO & TODO \\
\bottomrule
\end{tabular}
\caption{Illustrative subset of FormalAVS. Solve outcomes
(close vs open) per single attempt. Full benchmark numbers in
\citet{anonymous2026formalavs}.}
\label{tab:illustrative-comparison}
\end{table}

The wider table on the full benchmark subset is in
Appendix~\ref{sec:appendix-comparison}.

\subsection{Limitations}
\label{sec:discussion-limitations}

Three concrete gaps. (i) Domain scope: \texttt{pythia} is
statistics-domain only. It does not address general proof automation
in Lean 4. (ii) Toolchain pin: the library ships pinned to Lean
4.28 + Mathlib v4.28.0. Mathlib evolves quickly. Each major Mathlib
release typically requires a porting pass. (iii) Aristotle
dependency: the hardest closures route to a proprietary external
prover. The \texttt{z3\_check} oracle is open-source. The Aristotle
oracle is not. The library is usable without Aristotle for any goal
that closes via the \texttt{pythia} dispatcher chain or
\texttt{z3\_check}. A subset of structural Mathlib gaps requires
either Aristotle or a hand proof.


{\sloppy
\bibliographystyle{icml2026}
\bibliography{references}
}

\appendix

\section{Library inventory}
\label{sec:appendix-inventory}

The shipped library at the AI4MATH 2026 submission cut spans the
modules below. Each line lists module path and one-line summary.
All listed modules are axiom-clean against
$\mathcal{A}_{\text{core}} = \{\texttt{propext},
\texttt{Classical.choice}, \texttt{Quot.sound}\}$ unless explicitly
flagged as scaffold.

\subsection{Anytime-valid inference}
\begin{itemize}[leftmargin=1.3em,itemsep=0.2em]
\item \texttt{Pythia.HowardRamdasCS}: Howard-Ramdas Ville bound for
  self-normalised confidence sequences (rate
  $\eta_{\mathrm{HR}}(b) = \sqrt{b \log 2}$). Axiom-clean.
\item \texttt{Pythia.BettingCS}: betting-CS coverage certificate
  (rate $\eta_{\mathrm{betting}}(b) = 1/\sqrt{b\log 2 + 1}$).
  Axiom-clean.
\item \texttt{Pythia.BettingStrategy}: betting-strategy admissibility
  scaffolding for the \texttt{anytime\_valid} dispatch.
\item \texttt{Pythia.VilleSupermartingale}: countable-time and
  finite-horizon Ville inequalities. The
  \texttt{ville\_supermartingale}, \texttt{ville\_ineq}, and
  \texttt{ville\_supermartingale\_finite} lemmas are the workhorse.
\item \texttt{Pythia.EDetector}: e-process formalism, with
  \texttt{combine\_eprocesses\_avg} (average-merge) as the workhorse
  closer. Anytime-valid Type-I error control via Markov-at-stopping.
\end{itemize}

\subsection{Concentration}
\begin{itemize}[leftmargin=1.3em,itemsep=0.2em]
\item \texttt{Pythia.SubGaussianMG}: sub-Gaussian martingale
  concentration. The Chernoff exponent $\exp(-\tau^2/(2\sigma^2 N))$
  bound. Axiom-clean.
\item \texttt{Pythia.SubGamma}: sub-gamma tail bound. Axiom-clean.
\item \texttt{Pythia.MGFBoundedSubGamma}: bounded-RV
  Bernstein-Bennett MGF bound,
  $\mathbb{E}[\exp(\lambda X)]
   \leq \exp(\sigma^2 \lambda^2 / (2(1 - b\lambda/3)))$
  for centred $|X| \leq b$. Axiom-clean.
\item \texttt{Pythia.Bernstein}: iid and martingale Bernstein
  applications.
\end{itemize}

\subsection{Matrix concentration}
\begin{itemize}[leftmargin=1.3em,itemsep=0.2em]
\item \texttt{Pythia.MatrixLieb}: Lieb concavity step
  \citep{lieb1973} and operator-monotone log composition. Scaffold.
  Dependent operator-trace inequalities route to Aristotle.
\item \texttt{Pythia.MatrixBernstein}: iid matrix-Bernstein head
  theorem \citep{tropp2012usermatrix}. Scaffold.
\item \texttt{Pythia.MatrixBernsteinFull}: martingale extension.
  Scaffold.
\end{itemize}

\subsection{Other modules}
\begin{itemize}[leftmargin=1.3em,itemsep=0.2em]
\item The \texttt{Pythia.Control} subpackage carries
  \texttt{LyapunovDiscrete}, the discrete-time Lyapunov stability
  result for linear systems.
\item \texttt{Pythia.Queueing.ErlangB}: Erlang-B blocking formula.
\item \texttt{Pythia.Queueing.LittlesLaw}: Little's law steady-state
  identity.
\item \texttt{Pythia.TimeSeries.} \texttt{WoldDecomposition} \citep{wold1938}.
\item \texttt{Pythia.Risk.CVaR} \citep{rockafellaruryasev2000}: the
  Rockafellar-Uryasev variational characterisation of CVaR.
\item \texttt{Pythia.Asymptotics.DeltaMethod},
  \texttt{Pythia.Asymptotics.} \texttt{DeltaMethodMulti}: scalar and
  multivariate delta method.
\item \texttt{Pythia.InfoTheory.} \texttt{BretagnolleHuberBinary}: binary-case
  total-variation bound via the Bhattacharyya bridge.
\item The \texttt{Pythia.MeasureTheory} subpackage carries
  \texttt{OptionalStoppingUnbounded} and \texttt{PathMeasureRN}.
  Both are structural support for the Ville chain.
\end{itemize}

\section{Worked examples}
\label{sec:appendix-worked}

Each example below is a literal regression test from the
\texttt{pythia} test suite. Every \texttt{example} compiles
axiom-clean against $\mathcal{A}_{\text{core}}$.

\subsection{anytime\_valid: Ville-bound closer}
\label{sec:appendix-worked-av}

The countable-time variant closes the canonical Ville maximal
inequality on a non-negative supermartingale.

\begin{lstlisting}
import Pythia.Tactic.AnytimeValid
open MeasureTheory ProbabilityTheory ENNReal Pythia

example {Omega : Type*} {m0 : MeasurableSpace Omega}
    {mu : Measure Omega} [IsFiniteMeasure mu]
    {f : Nat -> Omega -> Real}
    {F : Filtration Nat m0}
    (hsup : Supermartingale f F mu)
    (hnn  : forall t omega, 0 <= f t omega)
    (hint : Integrable (f 0) mu)
    {c : Real} (hc : 0 < c) :
    mu {omega | exists t : Nat, f t omega >= c}
      <= (Integral.integral mu (f 0)).toNNReal / c.toNNReal := by
  anytime_valid
\end{lstlisting}

The tactic dispatches via the
\texttt{ville\_supermartingale} \texttt{exact} path. Hypotheses are
discharged by \texttt{assumption}. No registered rule is consulted.
The kernel-checked term reduces under
$\{\texttt{propext}, \texttt{Classical.choice}, \texttt{Quot.sound}\}$.

\subsection{z3\_check: linear-real-arithmetic with reconstruction}
\label{sec:appendix-worked-z3}

A three-step linear chain over $\mathbb{R}$. Z3 is queried as a
ranking filter. \texttt{linarith} produces the actual proof term.

\begin{lstlisting}
import Pythia.Tactic.Z3Check
open Pythia

example {a b c : Real} (h1 : a < b) (h2 : b <= c) :
    a < c := by
  z3_check
\end{lstlisting}

Trace: the encoder produces the SMT-LIB query
\texttt{(set-logic QF\_LRA) ... (assert (< v\_a v\_b)) (assert
(<= v\_b v\_c)) (assert (not (< v\_a v\_c))) (check-sat)}, which Z3
returns \texttt{unsat}. \texttt{linarith} then constructs the
Farkas certificate inside Lean. Z3 cannot smuggle axioms. If
\texttt{linarith} disagreed, the tactic would fail loud.

\subsection{pythia: chain-rule fall-through to aesop}
\label{sec:appendix-worked-pythia-fall}

A goal with no matching \texttt{@[stat\_lemma]} but a closed-form
chain in \texttt{linarith}. Stage 1 (aesop with the \texttt{Pythia}
ruleset) declines. Stage 2 (the cleanup chain combining
\texttt{simp}, \texttt{omega}, \texttt{linarith}, and
\texttt{positivity}) closes.

\begin{lstlisting}
import Pythia.Tactic.Pythia
open Pythia

example (a b c d : Real)
    (h1 : a <= b) (h2 : b <= c) (h3 : c <= d) :
    a <= d := by
  pythia
\end{lstlisting}

The dispatcher's \texttt{first}-combinator tries aesop, observes no
applicable \texttt{@[stat\_lemma]} closer, then tries the cleanup
chain. \texttt{linarith} closes the chain in one step.

\subsection{Registry-extend: tagging a user theorem}
\label{sec:appendix-worked-extend}

User extensions consist of tagging library lemmas with the
appropriate attribute. Below: a trivial real-arithmetic identity
registered for \texttt{pythia} via \texttt{@[stat\_lemma]}.

\begin{lstlisting}
import Pythia.Tactic.Pythia
open Pythia

@[stat_lemma]
theorem add_zero_real (x : Real) : x + 0 = x := by ring

-- pythia now closes goals matching the registered lemma:
example (x y : Real) : (x + 0) + (y + 0) = x + y := by
  pythia
\end{lstlisting}

The attribute re-elaborates as
\texttt{@[aesop safe apply (rule\_sets := [Pythia])]}. The lemma
then joins the \texttt{Pythia} aesop ruleset and is consulted by
\texttt{pythia} on every subsequent invocation. No fork or rebuild
of the tactic implementation is required.

\section{Reproduction}
\label{sec:appendix-repro}

\paragraph{Toolchain pin.}
Lean 4.28 + Mathlib v4.28.0. The repository's
\texttt{lean-toolchain} file declares the exact toolchain. The
Mathlib dependency in \texttt{lakefile.toml} pins the v4.28.0 release
tag.

\paragraph{Build.}
The library URL is redacted for double-blind review (full URL in
camera-ready). Reproduction commands once the URL is supplied:
\begin{verbatim}
git clone <pythia-url-redacted>
cd pythia
lake exe cache get
lake build
\end{verbatim}

\paragraph{Axiom audit.}
Each shipped module ends with a \texttt{\#print axioms} block. The
expected audit set is
$\{\texttt{propext}, \texttt{Classical.choice}, \texttt{Quot.sound}\}$.
Any deviation fails the audit.

\section{Related work, extended}
\label{sec:appendix-related}

\subsection{Lean 4 tactic infrastructure}
Mathlib ships a tactic suite the \texttt{pythia} dispatcher composes
with. \texttt{linarith} produces Farkas certificates and
\texttt{polyrith} routes polynomial-ideal goals via Sage with
kernel-checked reconstruction. \texttt{nlinarith} extends
\texttt{linarith} to a small non-linear fragment. Sign goals route
through \texttt{positivity} and $\sigma$-algebra-membership goals
through \texttt{measurability}. Generalised-congruence monotonicity
is discharged by \texttt{gcongr}. Monotone-composition goals route
through \texttt{bound}, which uses an attribute-driven rule
database.
The closest Mathlib relative to \texttt{pythia} is \texttt{bound}:
the \texttt{@[stats\_ineq]} attribute literally re-tags lemmas as
\texttt{@[bound]}, so the underlying \texttt{bound} dispatch picks
them up automatically. Aesop \citep{aesop2023} supplies the
\texttt{rule\_sets} mechanism the \texttt{pythia} dispatcher uses.
\texttt{pythia} is therefore a thin domain-registration layer over
\texttt{aesop} plus a Z3 oracle.

\subsection{Sledgehammer (Isabelle)}
Sledgehammer
\citep{paulsonblanchette2010threeyears, blanchettesledgehammer2011}
dispatches Isabelle/HOL goals to Vampire, E, Z3, CVC4, with proof
reconstruction inside Isabelle's kernel via Metis or
\texttt{smt}-tactic replay. The reconstruction discipline is the
template for both CoqHammer and \texttt{z3\_check}. There is no Lean
4 equivalent at the time of writing. \texttt{pythia} does not
attempt the Isabelle Sledgehammer's full premise-selection front
end. The library ships only the linear-real-arithmetic instance via
\texttt{z3\_check}.

\subsection{CoqHammer (Coq)}
CoqHammer \citep{czajkakaliszyk2018coqhammer} dispatches Coq goals
to Vampire, E, Z3, CVC4 with kernel-checked reconstruction. Premise
selection in CoqHammer is k-NN over a featurised representation of
the Coq context. The reconstruction step is a small certifying tactic
(\texttt{sauto}) that replays the ATP's witness inside the Coq
kernel. \texttt{z3\_check}'s reconstruction via \texttt{linarith}
mirrors this template at small scale. The external prover supplies
the verdict and the in-kernel tactic supplies the proof term. The
kernel is the final arbiter. We do not claim a full Lean-side reimplementation of
CoqHammer's premise-selection layer.

\subsection{Statistical formalisation libraries}
Mathlib's \texttt{Probability.Martingale} module supplies Doob
decomposition and Azuma's inequality but does not instantiate the
canonical confidence-sequence families, the per-family rate
functions, or Ville's inequality for the betting construction.
The \texttt{pythia} library closes that gap. Coq's
mathcomp.analysis and the Isabelle/HOL probability
formalisations cover overlapping but disjoint sub-areas, and
are not attribute-driven hammers.



\end{document}
